% v12 staging module: complete fixed-denominator factorization and resonances.

\section{固定した分母の partner と二つの無限共鳴}
\label{fp12:section}

二高分母の正局所模型で，一方の分母を固定したときの相手分母を調べる．
標準形の二つの高分母の役割は対称ではないため，固定した分母が
canonical な $A$ 側にある場合と $B$ 側にある場合を両方調べる必要がある．

\begin{theorem}[固定分母の有限な factorization と共鳴]
\label{fp12:partner-structure}
奇整数 $n>1$ に対し，正の整数局所模型 $(n,m,1,1)$ を持つ奇整数 $m>1$ の集合を
$\mathcal P(n)$ とする．この集合は $n\notin\{7,13\}$ なら有限である．
$n=7,13$ の場合にはそれぞれ
\[
 \{6k+1:k\geq1\}\subset\mathcal P(7),\qquad
 \{12k+1:k\geq1\}\subset\mathcal P(13)
\]
となり，これらの表示された無限集合を除く partners は有限個である．
従って $\mathcal P(n)$ が無限であることと $n\in\{7,13\}$ は同値である．
全ての有限な候補も，以下の二方向の正約数計算で上限なしに列挙できる．
\end{theorem}

\begin{proof}
正の二高分母標準形を
\[
 x=\frac{a-1}{r},\quad y=\frac{b-1}{r+1},\quad
 u=-s<0<v=z,\quad
 S=\frac{a-1}{s},\quad T=\frac{b-1}{z},
\]
\[
 A=a+x-s,\qquad B=b-y+z,\qquad AB=ab-1
\]
と書く．

\paragraph{固定分母が canonical な $A$ 側にある場合．}
$a=n$ とし，まず正約数 $r,s\mid n-1$ を選ぶ．
\[
 M=1+rs,\qquad C_A=(r+1)(s-x-1)
\]
と置けば，pair product は
\begin{equation}\label{fp12:A-factorization}
 z\{A(r+1)-MT\}=C_A
\end{equation}
と同値である．$C_A\ne0$ なら $z$ は $|C_A|$ の正約数なので，
\[
 T=\frac{A(r+1)-C_A/z}{M},\qquad b=zT+1
\]
を計算し，$T\in\mathbb Z_{>0}$，$r+1\mid zT$，$b$ が奇数という条件を検査すればよい．
この計算は有限である．

$C_A=0$ なら $s=x+1\geq2$ であり，$s\mid a-1=r(s-1)$ より $r=ts$，$t\geq1$ となる．
このとき
\[
 A=a-1=M-r-1,\qquad M=ts^2+1.
\]
式~\eqref{fp12:A-factorization} から $M\mid A(r+1)$，従って $M\mid(r+1)^2$ である．
しかし
\[
 (r+1)^2=tM+\{t(2s-1)+1\},\qquad
 0<t(2s-1)+1<M,
\]
最後の差が $t(s-1)^2>0$ なので矛盾する．
従ってこの向きに無限の枝はない．

\paragraph{固定分母が canonical な $B$ 側にある場合．}
$b=n$ とし，$r+1\mid n-1$，$r\geq1$ と正約数 $z\mid n-1$ を選ぶ．
\[
 N=1+(r+1)z,\qquad C_B=r(y-z-1)
\]
と置けば，pair product は
\begin{equation}\label{fp12:B-factorization}
 s\{NS-rB\}=C_B
\end{equation}
と同値である．$C_B\ne0$ なら $s$ は $|C_B|$ の正約数であり，
\[
 S=\frac{rB+C_B/s}{N},\qquad a=sS+1
\]
に対して $S\in\mathbb Z_{>0}$，$r\mid sS$，$a$ が奇数という条件を検査する．
この計算も有限である．

残る $C_B=0$ は $y=z+1$ と同値である．
$z\mid b-1=(r+1)(z+1)$ より $r+1=tz$，$t\geq1$ であり，$tz\geq2$ である．
\[
 b=tz(z+1)+1,\qquad N=tz^2+1,\qquad B=N+r.
\]
$S=rB/N$ の整数性は $N\mid r^2$，すなわち
\begin{equation}\label{fp12:resonance-divisibility}
 tz^2+1\mid t(2z+1)-1
\end{equation}
と同値である．$z\geq3$ なら右辺は正で左辺より小さいため不可能である．
$z=1$ なら $t\geq2$ であり，$t+1\mid3t-1$ は $t+1\mid4$ と同値なので $t=3$ となる．
$z=2$ なら $4t+1\mid5t-1$ である．$t=1$ は不可能であり，$t\geq2$ では
\[
 4t+1\leq5t-1<2(4t+1)
\]
なので $t=2$ に限られる．従って，共鳴する固定データはちょうど
\[
 (b,r,z,S,T)=(7,2,1,3,6),\qquad (13,3,2,4,6)
\]
である．

前者では $r\mid sS$ により $s=2k$，後者では $s=3k$ となる．
従って全ての共鳴模型は，それぞれ
\[
 \begin{array}{c|cccccc}
  b&a&r&s&z&S&T\\\hline
  7&6k+1&2&2k&1&3&6\\
  13&12k+1&3&3k&2&4&6
 \end{array}\qquad(k\geq1)
\]
にちょうど一致する．どちらも標準形の全整数条件を満たし，
$u=-s<0<v=z$ なので正の weights を持つ．
各固定分母の他の全枝は先の有限計算に含まれる．
これは $b=7,13$ の全ての局所模型がこの二族に入るという主張ではなく，
有限 factorization の右辺が零となる枝を全て分類したものである．
\end{proof}

\begin{proposition}[三と五の全 partners と負の低分母 entry]
\label{fp12:small-partners}
分母に上限を置かずに
\[
 \mathcal P(3)=\{7,11\},\qquad
 \mathcal P(5)=\{13,21,31,37\}
\]
である．これらの小さい分母は canonical な $A$ 側にしか現れず，
対応する標準形の全データは次の表で尽きる．
\begin{center}
\small
\begin{tabular}{ccccccc}
 $a$&$b$&$r$&$s$&$z$&$(S,T)$&$d$\\\hline
 3&7&1&1&1&$(2,6)$&$-5$\\
 3&11&1&1&2&$(2,5)$&$-4$\\
 5&13&1&1&1&$(4,12)$&$-9$\\
 5&21&1&1&2&$(4,10)$&$-7$\\
 5&31&1&2&6&$(2,5)$&$-4$\\
 5&37&1&1&4&$(4,9)$&$-6$
\end{tabular}
\end{center}
従って，小さい分母三の row から低分母へ向かう負の entry は常に $-1$ であり，
分母五の row では $-1$ または $-2$ に限られる．
\end{proposition}

\begin{proof}
定理~\ref{fp12:partner-structure} の有限計算を両方向で行う．
$a=3$ の $A$ 側では $r,s\in\{1,2\}$，$a=5$ では $r,s\in\{1,2,4\}$ を調べるだけである．
$T$ の正整数性と $r+1\mid zT$ を満たすものは，三の場合に
$(r,s,z)=(1,1,1),(1,1,2)$，五の場合に
$(1,1,1),(1,1,2),(1,1,4),(1,2,6)$ となる．
$B$ 側では，三の場合に $r+1=2$，五の場合に $r+1\in\{2,4\}$ であるが，
式~\eqref{fp12:B-factorization} の候補で正整数 $S$ と $r\mid sS$ を同時に満たすものはない．
表の全模型は正の weights を持つ．標準形の負の低分母 entry が $-s$ であることから，
最後の主張も従う．これは $b$ の探索上限を用いた列挙ではなく，
固定分母に対して完全であることが証明された正約数計算である．
\end{proof}
